\documentclass[10pt,a4paper]{amsart}
\usepackage[margin=19mm]{geometry}
\usepackage{amsmath,amssymb,mathtools}
\usepackage[scr=rsfs,cal=euler]{mathalfa}
\usepackage{xfrac}
\usepackage[unicode,hidelinks,bookmarks=false]{hyperref}
\newcommand{\ie}{i.e.\ }
\newcommand{\cf}{cf.\ }
\newcommand{\resp}{respectively\ }
\newcommand{\Subsection}{\subsection}
\providecommand{\nobreakdash}{\nobreak\hbox{-}}
\setlength{\emergencystretch}{2em}
\multlinegap0pt
\usepackage{mdframed}
\usepackage{mdframed}
\usepackage{mdframed}
\usepackage{mdframed}
\usepackage{amssymb}
\usepackage{dsfont}
\usepackage{mdframed}
\usepackage{enumerate}
\usepackage[normalem]{ulem}
\usepackage{tikz}
\usepackage{mathtools}
\newtheorem{theorem}{Theorem}
\newcommand{\RR}{{\mathbb R}}
\newcommand{\dens}{{\rm dens\,}}
\title{Bounded common fundamental domains for two lattices}
\author{Sigrid Grepstad, Mihail N. Kolountzakis}
\date{Source: arXiv:2507.00604v3 (CC BY 4.0)}
\begin{document}
\maketitle

\noindent\emph{Source and licence.} Bounded common fundamental domains for two lattices. Version arXiv:2507.00604v3 (CC BY 4.0), distributed by arXiv under the Creative Commons Attribution 4.0 International licence, http://creativecommons.org/licenses/by/4.0/, as recorded in the arXiv licence metadata for this exact version and shown on the abstract page https://arxiv.org/abs/2507.00604v3. Full licence notice: \texttt{licenses/arxiv-2507.00604-CC-BY-4.0.txt}.\par\smallskip

\noindent\textit{Editorial scope.} Counted unit: the article's theorem th:dense (source0.tex lines 277--279). The article's complete original proof of it is delivered with it (source0.tex lines 376--428, one \texttt{proof} environment of 53 lines, which the article places away from the statement). No mathematics is written, repaired or re-derived: the statement and the proof are the article's own lines, in the article's own notation, and no retained line is substituted or re-flowed.  Retained uncounted as the source-local context the proof needs: lines 354--358; lines 368--370.  Nothing was omitted from any retained span, so the package carries no omission record.  The retained lines cite 2 works, whose entries are transcribed verbatim from the article's own bibliography source in the e-print and delivered with short editorial labels recorded in the item's reference mapping.  No retained line contains a figure environment, an image-inclusion command or drawing code, so no image is delivered and the package declares no asset dependency.  Editorial formatting only, in addition: an AMS \texttt{amsart} preamble that restates the article's own declarations and macros used by the retained lines, namely the environments \texttt{theorem} on separate counters, together with the standard packages those lines need.  The retained environments print in this document as Theorem 1 in order of appearance, whereas the article prints the same items with the numbering of its own full document.  The complete native source of the article as distributed in the arXiv e-print for this exact version is bundled unchanged as \texttt{sources/180819/source0.tex}.
\begin{theorem}{\cite[Theorem 3.1]{DO1991}}\label{thm:do}
Let $\Gamma$ be a lattice in $\mathbb{R}^m \times \mathbb{R}^n$. If $W \subset \mathbb{R}^n$ is a parallelotope
\[W= \left\{ \sum_{k=1}^n t_k v_k \, : \, 0 \leq t_k < 1 \right\}\]
spanned by $n$ linearly independent vectors in $p_2(\Gamma)$, then the model set $\Lambda(\Gamma, W)$ is at bounded distance to a lattice in $\mathbb{R}^m$.
\end{theorem}

\begin{theorem}{\cite[Theorem 4.1]{EGKL2025}}\label{thm:equidecomp}
Let $\Gamma \subset \mathbb{R}^m \times \mathbb{R}^n$ be a lattice and let $W$ and $W'$ be bounded, measurable sets in $\mathbb{R}^n$ {where both $\partial W$ and $\partial W'$ have measure zero and $|W|=|W'|$}. If the model sets $\Lambda_W = \Lambda(\Gamma, W)$ and $\Lambda_{W'}=\Lambda(\Gamma, W')$ are bounded distance equivalent, then the windows $W$ and $W'$ are $p_2(\Gamma)$-equidecomposable up to measure zero.
\end{theorem}

\begin{theorem}\label{th:dense}
If $L, M \subseteq \RR^d$ are lattices of the same volume and $\overline{L+M}=\RR^d$ then there is a bounded, measurable $E \subseteq \RR^d$ such that $L \oplus E = M \oplus E = \RR^d$ are both tilings. 
\end{theorem}

\begin{proof}[Proof of Theorem \ref{th:dense}]
By abuse of notation let $L=L\mathbb{Z}^d$ and $M=M\mathbb{Z}^d$, where $L$ and $M$ are $d\times d$ non-singular matrices. Let $\Omega_L$ be the half-open parallelotope spanned by the columns of $L$ and $\Omega_M$ be the half-open parallelotope spanned by the columns of $M$. Then $\Omega_L$ and $\Omega_M$ are fundamental domains of the lattices $L$ and $M$, respectively. Since $L$ and $M$ are assumed to have equal volumes, we have $|\Omega_M| = |\Omega_L|$.

We now construct a lattice $\Gamma \subset \mathbb{R}^d \times \mathbb{R}^d$ (where $\Gamma$ again denotes both the lattice itself and its matrix representation) by letting
\begin{equation*}
\Gamma = 
\left(
\begin{array}{c c c c c}
\\
& & K & & \\
\\
 \hline
 \\
& L &  & M & \\ 
 \\
\end{array}
\right),
\end{equation*}
where $K$ is chosen to be a $d \times 2d$ matrix which acts as an injective map on $\mathbb{Z}^{2d}$. With the cut-and-projection construction in mind, note that this ensures that the projection $p_1$ is injective when restricted to the lattice $\Gamma$. This requirement for $K = [K_1 | K_2]$ is a generic condition on $K$ (say, the matrices $K$ failing it have measure $0$) hence the matrix $K_1$ can also be assumed to be nonsingular, in which case we can apply the block-determinant formula and obtain
\begin{equation}\label{block-det}
\det\Gamma = \det K_1 \det(M-L K_1^{-1} K_2).
\end{equation}
In order to guarantee the non-singularity of $\Gamma$ it is enough to pick the matrix $K_2$ to have very small elements. Since $M$ is nonsingular and the set of nonsingular matrices is open it follows that the second factor on the right-hand side of  \eqref{block-det} is also non-zero and thus $\Gamma$ is a proper lattice in $\RR^{2d}$. Finally, since $\overline{L+M}=\mathbb{R}^d$ by assumption, we know that $p_2(\Gamma)$ is dense in $\mathbb{R}^d$.


We now consider the two model sets 
\begin{equation*}
\Lambda_L= \Lambda(\Gamma, \Omega_L) = \left\{ p_1(\gamma) \, : \, \gamma \in \Gamma , p_2(\gamma) \in \Omega_L \right\}
\end{equation*}
and
\begin{equation*}
\Lambda_M = \Lambda(\Gamma, \Omega_M) = \left\{ p_1(\gamma) \, : \, \gamma \in \Gamma , p_2(\gamma) \in \Omega_M \right\}.
\end{equation*}
Since $p_2(\Gamma) = L + M$, we see that both $\Omega_L$ and $\Omega_M$ are windows spanned by $d$ linearly independent vectors in $p_2(\Gamma)$. Thus by Theorem \ref{thm:do}, both $\Lambda_L$ and $\Lambda_M$ are bounded distance equivalent to a lattice. Moreover, by assumption we have $|\Omega_L|=|\Omega_M|$, so $\dens \Lambda_L= \dens \Lambda_M$. This implies that the model sets $\Lambda_L$ and $\Lambda_M$ are bounded distance equivalent to lattices of equal density, and thus also {at bounded distance from each other}. We thus conclude from Theorem \ref{thm:equidecomp} that we can find a partition of $\Omega_L$ into subsets $S_1, \ldots , S_N$ and elements $\gamma_1 , \ldots , \gamma_N \in \Gamma$ such that 
\begin{equation}\label{eq:fdm}
\Omega_M = \bigcup_{i=1}^N \underbrace{\left(S_i + p_2(\gamma_i) \right)}_{S_i'} = \bigcup_{i=1}^N S_i',
\end{equation}
where we understand this equality to hold up to measure zero.

Finally, we observe that 
\begin{equation*}
p_2(\gamma_i) = \ell_i + m_i, 
\end{equation*}
for every $i=1, \ldots , N$, where $\ell_i \in L$ and $m_i \in M$. It follows that
\begin{equation*}
E=\bigcup_{i=1}^N (S_i' - m_i) = \bigcup_{i=1}^N (S_i + \ell_i)
\end{equation*}
is a fundamental domain for both $M$ and $L$ by \eqref{eq:fdm} and the fact that $(S_i)_{i=1}^N$ is a partition of $\Omega_L$. We thus have
\begin{equation*}
L \oplus E = M \oplus E = \mathbb{R}^d,
\end{equation*}
for a bounded measurable set $E \subset \RR^d$.
\end{proof}

\begin{thebibliography}{99}
\bibitem[DO1991]{DO1991} M.~Duneau and C.~Oguey. \newblock Bounded interpolations between lattices. \newblock {\em Journal of Physics A: Mathematical and General}, 24(2), 1991.
\bibitem[EGKL20]{EGKL2025} Mark~Mordechai Etkind, Sigrid Grepstad, Mihail~N. Kolountzakis, and Nir Lev. \newblock Bounded remainder sets, bounded distance equivalent cut-and-project sets, and equidecomposability. \newblock {\em arXiv preprint arXiv:2511.21148}, 2025.
\end{thebibliography}
\end{document}
